% !PRINTED 297
\weberpage{297}
%<<CONT
\weq{4}{
\begin{aligned}
\frac{d\varphi}{d\vartheta} &= -n R^n \sin n\vartheta - (n-1) p_1 R^{n-1} \sin[(n-1)\vartheta + q_1] \\
&\quad - \cdots - p_{n-1} \sin(\vartheta + q_{n-1}) \\[4pt]
\frac{d\psi}{d\vartheta} &= n R^n \cos n\vartheta + (n-1) p_1 R^{n-1} \cos[(n-1)\vartheta + q_1)] \\
&\quad + \cdots + p_{n-1} \cos(\vartheta + q_{n-1}).
\end{aligned}
} % [sic] stray closing parenthesis after q_1 in the cos-bracket of the second line, as printed

Wir theilen jetzt die Kreisperipherie mit dem Radius $R$ in
$4n$ Theile ein von der Winkelgrösse
\[
2\omega = 2\,\frac{\pi}{4n},
\]
indem wir bei $\vartheta = -\omega$ anfangen. Die Theilpunkte sind also
\[
-\omega,\ \omega,\ 3\omega,\ 5\omega,\ \dots (8n - 3)\omega,
\]
und die Intervalle mögen der Reihe nach mit $1, 2, 3 \dots 4n$
bezeichnet sein. Die Fig. 15 zeigt diese Eintheilung für $n = 3$.

\weberfigureimage{weber-b1-fig15.png}{0.54}

In den Intervallen
\[
1, 3, 5 \dots (4n - 1)
\]
ist $\cos n\vartheta$, absolut genommen, grösser als $1 : \sqrt{2}$, und
abwechselnd von positivem und negativem Vorzeichen. In den
Intervallen
\[
2, 4, 6 \dots 4n
\]
ist $\sin n\vartheta$ absolut grösser als $1 : \sqrt{2}$ und gleichfalls
abwechselnd von positivem und negativem Vorzeichen.

Man kann aber $R$ so gross annehmen, dass in den Intervallen
$1, 3, 5 \dots (4n - 1)$ das Vorzeichen von $\varphi$ und $\dfrac{d\psi}{d\vartheta}$ mit dem
von $\cos n\vartheta$, und in den Intervallen $2, 4, 6 \dots 4n$ das Vorzeichen
von $\psi$ und $\dfrac{d\varphi}{d\vartheta}$ mit dem von $\sin n\vartheta$ übereinstimmt.

Da also $\varphi$ in den Intervallen $2, 4, 6 \dots 4n$ sein Zeichen
ändert, so muss es in jedem dieser Intervalle durch Null gehen,
und da die Abgeleitete $\dfrac{d\varphi}{d\vartheta}$ ihr Zeichen nicht ändert, so kann
es nur \gesperrt{einmal} in jedem durch Null gehen.
